\section{Unit failure heads and safe retained completion arcs}
\label{sec:unit-failure-sources}

All reachability, convenience, protection, and failure sets in this
section are defined in the original graph $D$.
For a far target $y\in U_D(x)$, protection means
\[
 y\to x,\qquad N_D^-(y)\subseteq N_D^-(x)\setminus\{y\}.
\]
Write $\mathcal F(x)$ for the unprotected far targets of $x$, and
$t(x)=|\mathcal F(x)|$.
A protected target has the useful reachability implication
\[
 \begin{gathered}
 y\text{ protected for }x\\
 \Longrightarrow\quad
 [z\leadsto y\ (\le2)\ \Longrightarrow\ z\leadsto x\ (\le2)].
 \end{gathered}
\]
This follows directly from the incoming-neighborhood containment.

For a missing pair $ab$, let
\[
 F_{ab}=\{x:x\to a,\ b\in U_D(x)\}.
\]
The direction $a\to b$ is convenient precisely when $F_{ab}$ is empty.
The pair is bad precisely when both $F_{ab}$ and $F_{ba}$ are nonempty.
If it is bad and $x\in F_{ab}$, then $b\in\mathcal F(x)$:
an opposite-direction witness would force an arc toward $x$ if $b$
were protected, contradicting the cross-witness nonadjacency.

Conveniently orient all good missing pairs using the canonical rule:
choose the only convenient direction when unique; if both directions
are convenient, follow a fixed linear extension of the protected-target
strict partial order.  This rule has the compatibility property
\[
 \begin{gathered}
 y\text{ protected for }f,\quad f\to q\text{ canonically added good arc},\\
 qy\text{ missing}\quad\Longrightarrow\quad y\to q\text{ is chosen}.
 \end{gathered}
 \tag{UF1}\label{eq:unit-good-compatibility}
\]
In particular, no cardinality bound on the failure-source set is
needed here.  Indeed $G(f\to q)$ implies $G(y\to q)$, since every
predecessor of $y$ is a predecessor of $f$; hence $qy$ is automatically
good.  If its opposite direction is also convenient, then
$G(q\to y)$ implies $G(q\to f)$ by the reachability implication above.
Both pairs would then be oriented by the reference order, and the
choices $f\to q\to y$ would contradict $y$ preceding $f$.

\begin{lemma}[Safe retention with unit failure heads]
\label{lem:unit-safe-retained}
Suppose every directional failure source of a bad missing pair has
$t(x)\le1$.  Complete the good pairs canonically, and orient the bad
pairs arbitrarily to form a tournament $T$.
For any added arc $f\to q$, reverse all the other added out-arcs
of $f$ inward to form $T'$.
Then
\[
 N_{T'}^+(f)=N_D^+(f)\dot\cup\{q\},
\]
\[
 {
 N_{T'}^{++}(f)\subseteq
 \bigl(N_D^{++}(f)\setminus\{q\}\bigr)
 \cup\bigl(\mathcal F(f)\setminus\{q\}\bigr),}
 \tag{UF2}\label{eq:unit-retained-envelope}
\]
and consequently
\[
 {d_{T'}^{++}(f)\le d_D^{++}(f)+t(f)-1.}
\]
The retained pair $fq$ may be good or bad.
\end{lemma}
\begin{proof}
The first-neighborhood assertion is immediate.
Consider a second-neighborhood path $f\to a\to y$ in $T'$.
If $a$ is an original outneighbor of $f$, an original or convenient
second arc makes $y$ originally reachable within two steps.
A nonconvenient bad second arc can only introduce an unprotected
far head of $f$.  Second arcs in such a path are untouched by the
modification, since they do not involve $f$.

Now let $a=q$.  A protected far target $y$ of $f$ cannot receive
an original arc from $q$: that would force $q\to f$, whereas $fq$
is missing.  If $qy$ is missing and $fq$ is good, its chosen direction
is $y\to q$ by \eqref{eq:unit-good-compatibility}.

If $fq$ is bad, choose an opposite-direction failure source $r$:
$r\to q$ and $f\in U_D(r)$.
Protection of $y$ forces $y\in U_D(r)$ as well.
Thus the direction $q\to y$ fails at $r$.
If $qy$ were bad, its opposite witness would make
$y\in\mathcal F(r)$, while the bad pair $fq$ already makes
$f\in\mathcal F(r)$.  These two distinct heads contradict $t(r)\le1$.
Therefore $qy$ is good, with $y\to q$ its only convenient direction.
The arc $q\to y$ again cannot occur.

All new second heads are therefore in $\mathcal F(f)$, and $q$
itself is excluded because it is now a first neighbor.
The two sets in \eqref{eq:unit-retained-envelope} are disjoint.
Originally missing $fq$ puts $q$ in exactly one of
$N_D^{++}(f)$ and $U_D(f)$; in the latter case $q$ is necessarily
unprotected.  Excluding it removes one unit from their combined
cardinality, proving the final estimate.
\end{proof}

\subsection{A completion whose sink sources avoid failure}

Suppose all bad-edge directional failure sources lie in a set $Z$,
and all have $t\le1$.  Orient every missing pair inside $Z$ first,
using the canonical rule on good pairs.
Let $K$ consist of these oriented internal missing pairs and let
$S=\{z\in Z:d_K^+(z)=0\}$.
The set $S$ contains no missing pair.
Assume the already chosen bad internal directions have no failure
source in $S$.  Every remaining bad pair can then be oriented to
avoid failures in $S$: otherwise opposite witnesses in $S$ would
form a missing pair.  The remaining pairs do not change $K$.

If $t(z)\le g(z)+1$ for $z\in Z\setminus S$, this completion gives
a Seymour vertex in $D$.  To see this, choose a global median order
and its last vertex $f$.  If $f$ is outside $Z\setminus S$, reverse
all its added out-arcs inward; all its remaining first arcs are
original, and it is not a failure source of any chosen direction,
so its tournament SNP transfers directly to $D$.
Otherwise retain any added arc $f\to q$ in $K$, reversing the other
added out-arcs inward.  The order remains median.  By
Lemma~\ref{lem:unit-safe-retained}, its new deficiency is at least
$g(f)+2-t(f)\ge1$, contradicting the tournament feed theorem.
This reasoning does not infer SNP merely from acyclicity of a
dependency digraph.

\subsection{Four unit failure sources reduce to one reciprocal cycle}

\begin{proposition}\label{prop:four-unit-cycle-reduction}
Let $D$ be a finite nonempty all-deficient oriented graph.
Let $Z$ be the union of the directional failure-source sets of all
bad missing pairs. Suppose $|Z|\le4$ and $t(x)\le1$ for every $x\in Z$.
Then the only case not excluded by the preceding completion argument
has the following structure:
\begin{enumerate}
\item The four failure-source vertices induce a directed cycle of
length four, with the two diagonals missing.
\item These diagonals are the only bad missing pairs of $D$, and lose
to each other.
\item Every other missing pair is good.
\item Each failure-source vertex has exactly one unprotected far
head, its incoming neighbor on the four-cycle.
\end{enumerate}
\end{proposition}
\begin{proof}
Every bad-edge directional failure source is in $Z$.
An internal bad pair requires four distinct vertices of $Z$; if
$|Z|\le3$, the case of no internal bad pair below applies immediately.
Otherwise $|Z|=4$.
For each active such source, write $h(x)$ for its unique
unprotected far head.

Two bad missing pairs inside $Z$ cannot share an endpoint.
Suppose $ab$ and $ac$ were both bad, and let $d$ be the fourth vertex.
The witnesses for $ab$ must be $c,d$.
As $ac$ is missing, $c$ cannot point to $a$; hence $c\to b$,
$a\in U_D(c)$, and the opposite witness has $d\to a$ and
$b\in\mathcal F(d)$.
But $d$ is also a directional failure source for $ac$, which
would give it an unprotected head in $\{a,c\}$, different from $b$.
This contradicts $t(d)\le1$.
Thus internal bad pairs form a matching, with at most two members.

If there is no internal bad pair, use the canonical internal
directions.  If there is exactly one, say $ab$, its two witness
sources $c,d$ form a missing pair which is good.
Its canonical direction makes one of $c,d$, say $c$, a nonsink
of $K$.  Orient $ab$ in the direction failed by $c$.
There are no other possible directional failure sources for that
pair, because its endpoints are in $Z$ and only $c,d$ remain.
This internal direction therefore avoids every sink of $K$.
External bad pairs can be oriented to avoid those sinks as above.
The threshold $t\le1\le g+1$ applies to every nonsink source,
so both cases contradict that $D$ is all-deficient.

It remains that there are two internal bad pairs, say $ab$ and $cd$.
Label their witnesses so that $c\to a$, $d\to b$,
$h(c)=b$ and $h(d)=a$.
The witnesses for $cd$ are necessarily $a,b$.
Because opposite arcs are forbidden, they must satisfy
$a\to d$, $b\to c$, $h(a)=c$ and $h(b)=d$.
Thus the four cross pairs are precisely
$c\to a\to d\to b\to c$; the two diagonals are missing and
lose to each other.

Every bad pair has its two endpoints among the heads of its
opposite failure witnesses.  All four heads just identified
belong to $Z$, so there can be no other bad pair outside $Z$.
The internal matching classification leaves only $ab$ and $cd$.
The full dependency digraph consequently consists of this one
two-cycle and isolated missing pairs.
\end{proof}

\begin{theorem}[At most four unit failure sources]
\label{thm:four-unit-failure-sources}
Let $D$ be any finite nonempty oriented graph. If the union $Z$ of
the bad-edge directional failure sources has at most four vertices,
and $t(x)\le1$ for every $x\in Z$, then $D$ has a Seymour vertex.
\end{theorem}
\begin{proof}
Suppose instead that $D$ is all-deficient. By
Proposition~\ref{prop:four-unit-cycle-reduction}, only the reciprocal
cycle remains. Label its vertices $0,1,2,3$, with original arcs
$0\to1\to2\to3\to0$, missing diagonals $02,13$, and
\[
 h(0)=3,\quad h(1)=0,\quad h(2)=1,\quad h(3)=2.
\]
These diagonals are the only bad pairs. Complete them by $0\to2$
and $1\to3$, and complete all good pairs canonically, obtaining $T$.
The selected bad directions have failures only at $3$ and $0$,
respectively. Take a global median order $L$ of $T$ and let $f$ be
its last vertex.

If $f\notin\{0,3\}$, reverse all added outgoing arcs of $f$. The
result still has median order $L$, has the original first neighbourhood
of $f$, and has no new second head, because neither selected bad
direction fails at $f$. Its feed SNP transfers to $D$, a contradiction.
If $f=0$, retain the added arc $0\to2$ and reverse the other added
outgoing arcs. Lemma~\ref{lem:unit-safe-retained} gives deficiency
at least $g(0)+2-t(0)>0$, again contradicting the feed theorem.
Consequently every global median order of this fixed $T$ must end at $3$.

If $3$ has any added outgoing arc $3\to q$, retain it and reverse
the other added outgoing arcs. The same lemma gives positive
deficiency at this median feed, since $t(3)=1$. Thus $3$ has no
added outgoing arc at all, and $N_T^+(3)=N_D^+(3)$.
The only possible new second head at $3$ is its unique unprotected
far head $h(3)=2$. Therefore
\[
 N_T^{++}(3)\subseteq N_D^{++}(3)\cup\{2\}.
\]
The feed SNP and $g(3)>0$ imply
\[
 g(3)=1,\qquad d_T^+(3)=d_T^{++}(3).
\]

Use the strong form of the Havet--Thomasse feed theorem
\cite[Theorem~1 and its proof]{HavetThomasse2000}:
\[
 |N_T^+(3)|\le |G_L|\le |N_T^{++}(3)|,
\]
where $G_L$ consists of the incoming neighbours $y$ of $3$ for which
some outgoing neighbour $a$ occurs before $y$ in $L$ and satisfies
$a\to y$. The equality just obtained forces
$|N_T^+(3)|=|G_L|$.
The sedimentation lemma \cite[Lemma~1]{HavetThomasse2000} then states
that the order
\[
 \operatorname{Sed}(L)=
 (\text{incoming vertices outside }G_L,\ 3,\ N_T^+(3)\cup G_L)
\]
is still a global median order of the same tournament $T$, with each
displayed group preserving its old internal order. Its last vertex
is not $3$, since $N_T^+(3)$ is nonempty in an all-deficient graph.
This contradicts the fact that every global median order of $T$
ends at $3$. The reciprocal case is impossible as well.
\end{proof}

\begin{corollary}[No four-unit residual distribution]
\label{cor:no-four-unit-residuals}
In a strongly connected all-deficient graph minimal under arbitrary
nonempty arc-set deletion, residual four cannot be distributed as
$(1,1,1,1)$. Combined with
Proposition~\ref{prop:one-two-unit-source}, this supplies an intermediate exclusion; all remaining residual-four allocations are excluded in the residual-exclusion part of this supplement.
\end{corollary}
\begin{proof}
All bad-edge failure sources belong to the positive-residual support,
and residual one gives $t\le1$. Apply
Theorem~\ref{thm:four-unit-failure-sources}.
The support contains a missing pair, so a single positive vertex is
impossible. The remaining integer partitions are as stated.
\end{proof}

The sedimentation step uses a global median order and the equality
$|N_T^+(f)|=|G_L|$. It does not infer the existence of two feed vertices
from the existence of two Seymour vertices in a tournament, nor does
it assert that an arbitrary tournament admits two median feeds.
